Eine reale Stromquelle kann durch eine ideale Stromquelle mit dem Quellstrom $I_C$
und einem parallel geschalteten Innenwiderstand $R_I$ modelliert werden. An diese
Quelle ist ein Lastwiderstand $R_L$ parallel angeschlossen. Ziel ist es, die im Last-
widerstand umgesetzte Leistung in Abhängigkeit des Lastwiderstands zu untersuchen.
\begin{figure}[H]
    \centering
    \includegraphics[width=1\textwidth]{Stromquelle.png}
    \caption{Leistungsanpassung mit Stromquelle}
    \label{fig:Stromquelle}
\end{figure}

\noindent
Die Werte der betrachteten Stromquelle betragen:
\begin{equation*}
I_0 = 3\,\mathrm{A}
\end{equation*}
\begin{equation*}
R_I = 5\,\Omega
\end{equation*}

\noindent
Aufgrund der Parallelschaltung bietet es sich an, die Berechnungen mithilfe von Leit-
werten durchzuführen. Dazu werden der Innenleitwert sowie der Lastleitwert definiert
als:
\begin{equation*}
G_I = \frac{1}{R_I}
\end{equation*}
\begin{equation*}
G_L = \frac{1}{R_L}
\end{equation*}